% Chapter 1 -- Introduction to InitechDOS.
% Sources (Law 1): ADR-0003 DEC-01/02/06/07/09/11/12/13; os/milton/command.c
% (internal-command table, run_external PATH search .COM/.EXE/.BAT);
% os/milton/devices.c (CON PRN AUX CLOCK$ NUL); README.md (requirements);
% spec/hardware.json (cpu row); spec/dos_banner.txt.
\Gates{test-boot, test-shell, test-command, test-devices, test-sysinit, test-autoexec, test-hardware-spec}
\chapter{Introduction to InitechDOS}
\begin{ChapterCovers}
\item What an operating system does
\item The parts of InitechDOS
\item What you need to use InitechDOS
\item Drives, files and file names
\item Directories
\item Commands and messages
\end{ChapterCovers}

\section{What an Operating System Does}
\index{operating system}
A computer by itself does not know how to read a disk, show text on the
display, or run a program.  An \emph{operating system} is the program that
does these things for you and for the other programs you use.  It is the
first program the computer loads when you turn it on, and it stays in
control until you turn the computer off.

InitechDOS\index{InitechDOS} is the operating system for Initech business
workstations.  It does three kinds of work:
\begin{itemize}
\item It keeps track of the information stored on your disks, in units
  called \emph{files}.
\item It runs programs: word processors, spreadsheets, data managers and
  the programs your company writes for itself.
\item It accepts commands that you type at the keyboard and carries them
  out.
\end{itemize}
Most of the time you will be working in an application program.  You use
InitechDOS commands to prepare disks and files for that work and to look
after them afterwards: to see what files are on a disk, to make copies of
them, to remove files you no longer need, and to keep related files
together.

\section{The Parts of InitechDOS}
InitechDOS is made of three parts that work together.

\SideHead{Input/output system}The input/output system controls the devices
attached to the computer: the keyboard, the display, the disks, the printer,
the serial port and the clock.  When the computer starts, it also reads the
configuration file \file{CONFIG.SYS}\index{CONFIG.SYS@\texttt{CONFIG.SYS}},
which describes how your system is set up.

\SideHead{Kernel}The kernel manages files, directories and memory, and
loads and runs programs.  Programs ask the kernel for these services; you
do not deal with it directly.
\draft{ADR-0003 DEC-01 names the system files IO.SYS and INITDOS.SYS;
the shipped startup disk does not yet carry them as files, so this book does
not name them}

\SideHead{Command processor}The command processor, \file{COMMAND.COM}%
\index{COMMAND.COM@\texttt{COMMAND.COM}}, is the part you talk to.  It
displays the command prompt, reads each command you type, and either
carries the command out itself or finds the program you named and runs it.
When it starts, \file{COMMAND.COM} carries out the commands in the file
\file{AUTOEXEC.BAT}\index{AUTOEXEC.BAT@\texttt{AUTOEXEC.BAT}} if that file is
on drive~A.  Chapter~2 shows the command processor at work.

\section{What You Need}
\index{system requirements}
To use InitechDOS Version~3.30 you need:
\begin{itemize}
\item An Initech workstation or a compatible personal computer with an
  Intel 80486 or 80386 microprocessor
\item 4~MB of memory (2~MB minimum)
\item A fixed-disk controller and one fixed disk
\item One 5.25-inch or 3.5-inch diskette drive
\item A VGA-compatible display adapter and monitor
\end{itemize}
\draft{list from README.md; spec/hardware.json makes the 80486 the target and
the 80386 desirable -- confirm the wording and the memory figures with the
hardware contract before print}

\section{Drives and Disks}
\index{drives}\index{disks}
InitechDOS keeps files on disks.  Each disk drive is known by a letter
followed by a colon: the first drive is drive~A, written \file{A:}.
The letter of the drive you are working with appears in the command prompt,
\file{A:\textbackslash>}, which you will see in Chapter~2.

In this version of InitechDOS, the disk that holds your files is drive~A.
When you need to name a drive in a command, type its letter and a colon.
\draft{single drive: the shipped build mounts one FAT volume as drive A; a
second letter such as B: is answered "Bad command or file name" (screen
"errors"); the Reference will cover other drives when they are built}

\SideHead{Volume label}A disk may be given a name of up to eleven characters,
called its \emph{volume label}\index{volume label}.  InitechDOS shows the
label at the top of every directory listing, so that you can tell your
disks apart.  The practice diskette used in Chapter~2 has the label
\file{INITECH}.

\SideHead{Disk space}InitechDOS reports how much room is left on a disk at
the end of every directory listing, as a number of bytes.  A byte holds one
character of text.  A 3.5-inch high-capacity diskette holds 1,457,664 bytes
of files when it is empty.\index{bytes free}

\begin{Caution}
Do not remove a diskette from its drive while the drive light is on.  The
information on the diskette can be damaged.
\end{Caution}

\section{Files and File Names}
\index{files}\index{file names}
Everything InitechDOS keeps on a disk is kept in a \emph{file}: a letter, a
list of telephone numbers, a spreadsheet, or a program.  Each file on a disk
has a name, and no two files in the same directory can have the same name.

A file name has two parts: the \emph{name} and an optional
\emph{extension}\index{extensions}.
\begin{itemize}
\item The name is from one to eight characters long.
\item The extension is from one to three characters long.  It is separated
  from the name by a period.
\end{itemize}
For example, in the file name \file{MEMO.TXT} the name is \file{MEMO} and the
extension is \file{TXT}.  Use letters and numbers in your file names; they
are always safe.  You may type a file name in capital or small letters:
InitechDOS stores and displays it in capitals.

The extension usually tells you what kind of file it is.  Some extensions
have a meaning to InitechDOS:

\SideHead{.COM and .EXE}Program files.  To run the program, type its name
without the extension.

\SideHead{.BAT}Batch files: lists of InitechDOS commands that the command
processor carries out one after another.

\SideHead{.SYS}System files, such as \file{CONFIG.SYS}.

Other extensions, such as \file{.TXT} for text, are a matter of habit.
Choose extensions that help you recognize your files.

\SideHead{Device names}Some names are reserved for the devices attached to
the computer.  Do not use them as file names, with or without an
extension.\index{device names}
\par\medskip
\begin{tabular}{@{}l@{\quad}p{80mm}@{}}
{\imSideHeadFont Name} & {\imSideHeadFont Device}\\[2pt]
\cmd{CON}    & The console: the keyboard for input, the display for output\\
\cmd{PRN}    & The printer\\
\cmd{AUX}    & The auxiliary (serial) port\\
\cmd{CLOCK\$} & The system clock\\
\cmd{NUL}    & A device that discards whatever is written to it\\
\end{tabular}
\index{CON@\texttt{CON}}\index{PRN@\texttt{PRN}}\index{AUX@\texttt{AUX}}%
\index{NUL@\texttt{NUL}}

\SideHead{Wildcards}Some commands can act on a group of files at once.  To
name the group, use the \emph{wildcard}\index{wildcards} characters in the
file name.  An asterisk (\type{*}) stands for any characters from its
position to the end of the name or the extension; a question mark
(\type{?}) stands for any one character.  \file{*.TXT} means every file
with the extension \file{TXT}; \file{*.*} means every file.  Chapter~2 shows
which of the commands described there accept wildcards.

\section{Directories}
\index{directories}
A disk may hold a great many files.  To keep them in order, InitechDOS lets
you group files into \emph{directories}.  Every disk has one main
directory, the \emph{root directory}\index{root directory}, which is
created when the disk is formatted.  You can create further directories,
called \emph{subdirectories}\index{subdirectories}, inside the root
directory and inside other subdirectories, and give each one a name, in
the same form as a file name.

The directory you are working in is the \emph{current
directory}\index{current directory}.  The command prompt shows it: at
\file{A:\textbackslash>} you are in the root directory of drive~A, and at
\file{A:\textbackslash REPORTS>} you are in the subdirectory
\file{REPORTS}.

To name a file in another directory, write its \emph{path}\index{paths}:
the names of the directories that lead to it, each followed by a backslash
(\type{\textbackslash}).  A path that begins with a backslash starts at the
root directory.  \file{A:\textbackslash REPORTS\textbackslash MEMO.TXT} is
the file \file{MEMO.TXT} in the directory \file{REPORTS}, which is in the
root directory of drive~A.  Chapter~2 shows you how to create directories,
move between them, and remove them.

\section{Commands}
\index{commands}
You tell InitechDOS what to do by typing a \emph{command} at the command
prompt.  A command is a command name, sometimes followed by
\emph{parameters}\index{parameters} that tell the command what to work on;
for example, \type{TYPE MEMO.TXT}.  Separate the command name and each
parameter with a space.

\SideHead{Internal commands}Some commands are built into the command
processor and are always available.  They are called \emph{internal
commands}\index{internal commands}:
\par\medskip
\begin{tabular}{@{}l@{\quad}l@{\quad}l@{\quad}l@{}}
\cmd{BREAK} & \cmd{COPY} & \cmd{ECHO} & \cmd{REN} (\cmd{RENAME})\\
\cmd{CD} (\cmd{CHDIR}) & \cmd{DATE} & \cmd{EXIT} & \cmd{SET}\\
\cmd{CLS} & \cmd{DEL} (\cmd{ERASE}) & \cmd{MD} (\cmd{MKDIR}) & \cmd{TIME}\\
 & \cmd{DIR} & \cmd{PROMPT} & \cmd{TYPE}\\
 & & \cmd{RD} (\cmd{RMDIR}) & \cmd{VER}\\
\end{tabular}
\par\medskip
A name in parentheses is another name for the same command.

\SideHead{External commands}Any other command name is taken to be the name of
a program or batch file on disk; these are called \emph{external
commands}\index{external commands}.  The command processor looks for a file
with that name and the extension \file{.COM}, \file{.EXE} or \file{.BAT}, in
that order, first in the current directory and then in each directory
listed in the \cmd{PATH} environment variable.  If it finds none, it
displays:
\begin{TypeIn}
Bad command or file name
\end{TypeIn}
\index{Bad command or file name@\texttt{Bad command or file name}}
The same message appears when you mistype a command name.  Check the
spelling and type the command again.

\section{Messages}
\index{messages}
When InitechDOS cannot do what you asked, it tells you why in a short
message, such as \texttt{File not found}.  Each message always has exactly
the same wording, so that you can look it up.  Appendix~A lists the
messages; the \doctitle{InitechDOS Reference} explains each one and what
to do about it.
